% Source: structured_predictions.tex
\begin{notechapter}
\chapter{Structured Prediction}\label{ch:structured-prediction}

\section{Introduction}
Structured prediction is a subfield of machine learning where the goal 
is to predict structured outputs, such as sequences, trees, graphs, or 
other complex data structures, rather than simple scalar values or 
class labels. It is widely used in tasks where the output has inherent 
structure and dependencies, such as natural language processing, 
computer vision, and bioinformatics.

\section{What Is Structured Prediction?}
Structured prediction involves predicting outputs that are structured 
objects. Unlike traditional classification or regression tasks, where 
the output is a single value or label, structured prediction deals 
with outputs that have internal dependencies and relationships. 
Examples include:
\begin{itemize}
  \item \textbf{Sequences}: E.g., part-of-speech tagging, named entity 
    recognition, machine translation.
  \item \textbf{Trees}: E.g., syntactic parsing in natural language 
    processing.
  \item \textbf{Graphs}: E.g., predicting molecular structures in 
    chemistry.
  \item \textbf{Grids or Images}: E.g., semantic segmentation, image 
    captioning.
\end{itemize}

\section{Applications of Structured Prediction}
Structured prediction is used in a variety of domains:
\begin{itemize}
  \item \textbf{Natural Language Processing (NLP)}:
    \begin{itemize}
      \item \textbf{Sequence Labeling}: Predicting a sequence of labels 
        for a sequence of inputs (e.g., part-of-speech tagging, named 
        entity recognition).
      \item \textbf{Machine Translation}: Predicting a sequence of 
        words in a target language given a sequence of words in a 
        source language.
      \item \textbf{Parsing}: Predicting a syntactic tree for a 
        sentence.
    \end{itemize}
  \item \textbf{Computer Vision}:
    \begin{itemize}
      \item \textbf{Semantic Segmentation}: Assigning a label to each 
        pixel in an image.
      \item \textbf{Object Detection}: Predicting bounding boxes and 
        labels for objects in an image.
      \item \textbf{Image Captioning}: Generating a sequence of words 
        describing an image.
    \end{itemize}
  \item \textbf{Bioinformatics}:
    \begin{itemize}
      \item \textbf{Protein Structure Prediction}: Predicting the 3D 
        structure of a protein from its amino acid sequence.
      \item \textbf{Gene Prediction}: Identifying regions of DNA that 
        encode genes.
    \end{itemize}
\end{itemize}

\section{Approaches to Structured Prediction}
Several approaches are used to tackle structured prediction tasks:
\begin{itemize}
  \item \textbf{Probabilistic Models}:
    \begin{itemize}
      \item Use probabilistic frameworks to model dependencies between 
        output variables.
      \item Examples:
        \begin{itemize}
          \item \textbf{Hidden Markov Models (HMMs)}: For sequence 
            labeling tasks.
          \item \textbf{Conditional Random Fields (CRFs)}: A 
            discriminative model for structured prediction.
        \end{itemize}
    \end{itemize}
  \item \textbf{Neural Networks}:
    \begin{itemize}
      \item Use neural networks to model complex dependencies in 
        structured outputs.
      \item Examples:
        \begin{itemize}
          \item \textbf{Recurrent Neural Networks (RNNs)}: For sequence 
            prediction tasks.
          \item \textbf{Graph Neural Networks (GNNs)}: For 
            graph-structured data.
          \item \textbf{Transformers}: For tasks like machine 
            translation and text generation.
        \end{itemize}
    \end{itemize}
  \item \textbf{Structured Support Vector Machines (SSVMs)}:
    \begin{itemize}
      \item Extend support vector machines to handle structured outputs 
        by incorporating a loss function that accounts for the 
        structure of the output.
    \end{itemize}
  \item \textbf{Energy-Based Models}:
    \begin{itemize}
      \item Define an energy function that scores the compatibility 
        between inputs and structured outputs, then minimize the energy 
        during inference.
    \end{itemize}
\end{itemize}

\section{Challenges in Structured Prediction}
Structured prediction poses several challenges:
\begin{itemize}
  \item \textbf{Modeling Dependencies}:
    \begin{itemize}
      \item Capturing dependencies between elements of the output 
        structure (e.g., correlations between labels in a sequence).
    \end{itemize}
  \item \textbf{Inference}:
    \begin{itemize}
      \item Finding the best output structure given the input (e.g., 
        using dynamic programming, beam search, or other optimization 
        techniques).
    \end{itemize}
  \item \textbf{Training}:
    \begin{itemize}
      \item Learning model parameters efficiently, often requiring 
        specialized algorithms like backpropagation through time 
        (BPTT) for RNNs or message passing for CRFs.
    \end{itemize}
\end{itemize}

\section{Example: Sequence Labeling with CRFs}
In sequence labeling tasks like named entity recognition, a 
\textbf{Conditional Random Field (CRF)} can be used to model 
dependencies between labels. For example:
\begin{itemize}
  \item Input: A sequence of words \( \mathbf{x} = (x_1, x_2, \dots, 
    x_n) \).
  \item Output: A sequence of labels \( \mathbf{y} = (y_1, y_2, \dots, 
    y_n) \).
  \item The CRF models the conditional probability \( P(\mathbf{y} 
    \mid \mathbf{x}) \) and uses features that capture dependencies 
    between neighboring labels.
\end{itemize}

\section{Conclusion}
Structured prediction is a powerful framework for handling complex 
outputs with inherent structure and dependencies. It is widely used in 
tasks like natural language processing, computer vision, and 
bioinformatics. By leveraging probabilistic models, neural networks, 
and other approaches, structured prediction enables machines to solve 
challenging real-world problems.
\end{notechapter}
